\title{Direct Preference Optimization: Your Language Model is Secretly a Reward Model}

\begin{abstract}
Although large-scale unsupervised language models (LMs) acquire extensive general knowledge and some capabilities for reasoning, their behaviors are inherently challenging to shape because their training relies entirely on unsupervised methods. Common approaches to increase their controllability involve gathering human feedback on the relative quality of the LM outputs, subsequently refining the LM through supervised fine-tuning in accordance with these preferences—most notably through reinforcement learning from human feedback (RLHF). RLHF, however, is an intricate and often unstable process, consisting first of training a reward model that represents human judgments, followed by reinforcement learning on the large LM to maximize this estimated reward, all while attempting to prevent significant divergence from the pretrained parameters. In this work, we propose a novel reward model parameterization for RLHF, which permits direct derivation of the optimal policy in closed form, thereby reducing the conventional RLHF objective to a straightforward classification loss. This yields a new method, termed \textit{Direct Preference Optimization} (DPO), which is efficient, robust, and computationally streamlined: it avoids the need for sampling from the LM during fine-tuning and dispenses with extensive hyperparameter search. Through our experimental evaluation, we demonstrate that DPO is capable of fine-tuning LMs for alignment with human preference at least as effectively as prevailing methods. Importantly, DPO not only surpasses RLHF with PPO in managing sentiment control but also achieves parity or improvement in summarization and single-turn dialogue quality, all while substantially reducing complexity in both implementation and training.
\end{abstract}

\section{Introduction}
Large unsupervised language models (LMs), when trained on massive text corpora, often display remarkable emergent capabilities~\citep{chowdhery2022palm, brown2020language, touvron2023llama,bubeck2023sparks}. Nonetheless, these models are exposed to data produced by humans whose objectives, motivations, and expertise vary widely. Not all of these human-generated goals and skills are appropriate for imitation by AI; for instance, we may desire that an AI code assistant \textit{recognize} frequent programming errors in order to fix them, but we prefer that the assistant's generated code consistently reflect the exemplary, albeit less common, programming proficiency present in the training dataset. By the same token, we may want a language model to be \textit{informed} about a widespread misconception held by half the population, while we decidedly do not wish for it to affirm that misconception in half of its answers to relevant questions. Thus, it is essential to carefully \emph{select} the model's \emph{preferred outputs and conduct} from the broad array of its \textit{acquired knowledge and skills} in order to develop AI systems that are reliable, effective, and amenable to control \citep{ouyang2022training}. Most existing techniques attempt to align LMs with human values and preferences by leveraging reinforcement learning (RL), but we demonstrate that the standard RL-based objective can in fact be solved precisely with a straightforward binary cross-entropy loss, thus greatly streamlining the process of preference learning.

Broadly speaking, prevalent approaches infuse language models with desired behaviors via datasets that capture human preferences for responses considered safe and beneficial. This phase of preference alignment follows an initial round of unsupervised pre-training using extensive textual data. Although the most intuitive method for preference alignment is to apply supervised fine-tuning with human-provided examples of preferred outputs, the leading suite of approaches centers on reinforcement learning from human or AI feedback (RLHF/RLAIF; \citep{christiano2017deep,bai2022constitutional}). In RLHF, a reward model is first trained on collected human preference data, and then this reward model guides reinforcement learning to adjust the language model policy so that it produces outputs rated highly by the reward model, while restricting deviation from the base model. Despite the high-quality conversational and coding performances RLHF models can achieve, the RLHF process is considerably more intricate compared to supervised approaches, requiring the simultaneous training of multiple LMs and on-policy sampling during optimization, which leads to substantial computational resource demands.

This work presents a method for optimizing language models to better reflect human preferences, bypassing the need for explicit reward modeling or reinforcement learning techniques. We introduce \textit{{Direct Preference Optimization} (DPO)}, a procedure that achieves an objective equivalent to that of standard RLHF approaches (specifically, reward maximization with a KL-divergence penalty), yet stands out for its simplicity in both implementation and training. Conceptually, the DPO approach enhances the ratio of the log likelihood assigned to preferred responses over dispreferred ones, integrating a sample-specific importance weighting mechanism to mitigate model collapse, which we observe when employing a straightforward probability ratio objective. Similar to existing techniques, DPO builds on a theoretical preference framework—such as the Bradley-Terry model \cite{bradley1952rankanalysis}—to quantify how a reward function corresponds to collected human preference judgments. The distinguishing feature of DPO, however, lies in its use of a variable transformation that allows the preference loss to depend directly on the policy parameters, as opposed to first training a reward model. Consequently, when supplied with a dataset containing human choices among candidate outputs, DPO can train the policy with a simple binary cross entropy loss, yielding an optimal policy with respect to a reward function implicitly defined by the observed preferences.

The principal innovation of this work is the Direct Preference Optimization (DPO) algorithm, offering a reward-model-free alternative to reinforcement learning for preference-based language model training. Empirical results demonstrate that DPO matches or surpasses existing approaches, including those relying on PPO for RLHF, in applications such as sentiment control, summarization, and dialogue across language models scaling up to 6B parameters.

\section{Related Work}

Language models trained via self-supervision at increasing scales have demonstrated the capacity to solve certain tasks in a zero-shot manner \citep{radford2019language} or when provided with minimal prompt examples \citep{gpt3,megatron,chowdhery2022palm}. Nevertheless, their effectiveness on downstream benchmarks and their alignment with user intentions see substantial gains when further refined on corpora comprising instructions and human-generated completions \citep{mishra-etal-2022-cross,sanh2022multitask,chung2022scaling,thoppilan2022lamda}. This process, known as `instruction-tuning,' enhances LLMs' ability to generalize to instructions that fall outside their fine-tuning distribution and tends to increase their practical utility \citep{chung2022scaling}. Although instruction-tuning has proven fruitful, it is frequently easier to collect \textit{relative} human assessments of response quality than to obtain high-quality expert demonstrations. As a result, subsequent studies have trained LLMs on datasets capturing human preferences, resulting in advances in areas such as translation \citep{kreutzer-etal-2018-reliability}, summarization \citep{stiennon2022learning,ziegler2020finetuning}, story generation \citep{ziegler2020finetuning}, and instruction-following \citep{ouyang2022training,ramamurthy2023is}. In these approaches, a neural network reward function is trained to represent the observed preference data using a framework like the Bradley-Terry model \citep{bradley1952rankanalysis}. The language model is then fine-tuned by optimizing the reward signal, employing reinforcement learning techniques such as REINFORCE \citep{williams1992reinforce}, proximal policy optimization (PPO; \cite{schulman2017proximal}), or similar variants \citep{ramamurthy2023is}. A parallel research avenue utilizes LLMs that have already been refined for instruction following through human feedback to generate further synthetic preference data focusing on attributes like safety or harmlessness \citep{bai2022constitutional}. This is achieved using minimal human guidance, typically provided through textual rubrics that guide the model’s evaluation process. These approaches illustrate the intersection of two previously distinct research traditions: reinforcement learning applied to language model training across diverse objectives~\citep{Ranzato2015SequenceLT,paulus2018a,wu2018learning}, and the development of general strategies for leveraging human preferences in learning systems \citep{christiano2017deep,kupcsik2018learning}. Although leveraging relative human preferences has strong appeal, the fine-tuning of large language models with reinforcement learning poses significant practical hurdles. In response, the current work proposes a principled alternative for optimizing with relative preferences, without relying on reinforcement learning.

Beyond language applications, the study of policy learning from preferences has been explored within both bandit frameworks and reinforcement learning paradigms, yielding various methodological developments. The use of preferences or action rankings instead of explicit reward signals in contextual bandit problems is referred to as the contextual dueling bandit (CDB) model \citep{yue2012karmed,dudik2015contextual}. In these scenarios, where absolute reward values are unavailable, theoretical work on CDBs replaces the standard concept of an optimal policy with that of a \textit{von Neumann winner}—a policy whose average probability of outperforming any other policy reaches at least 50\% \citep{dudik2015contextual}. Notably, preference information in CDBs is typically collected in an online manner, whereas learning from human preferences generally entails training on a predetermined dataset of offline preference-labeled action pairs \citep{yan2022human}. 

Likewise, \textit{preference-based RL} (PbRL) focuses on leveraging binary preferences generated through an \textit{unknown} 'scoring' mechanism, replacing direct access to reward feedback \citep{BusaFekete2014,ruiz2023dueling}. Numerous PbRL algorithms have been proposed, some allowing the reuse of off-policy preference data, yet the common approach is to first model the latent scoring (reward) function before optimizing the policy using that model \citep{jain2013learning,BusaFekete2014,christiano2017deep,sadigh2017active,kupcsik2018learning}. In contrast to these methods, we introduce a unified policy learning framework that seeks to directly align policy outputs with observed preferences in a single stage.

\section{Preliminaries}\label{section:prelims}

This section reviews the RLHF procedure outlined by \citeauthor{ziegler2020finetuning}, subsequently developed in works such as \citep{stiennon2022learning, bai2022training, ouyang2022training}. The RLHF pipeline generally comprises three primary steps: (1) supervised fine-tuning (SFT); (2) collecting preference data and constructing a reward model; and (3) reinforcement learning for policy optimization.

\textbf{SFT}: The standard starting point for RLHF is to apply supervised learning, fine-tuning a pre-trained language model (LM) on a curated corpus for relevant downstream tasks (e.g., dialogue systems, summarization), resulting in the SFT policy $\pi^\text{SFT}$.

\textbf{Reward Modelling Phase}: Following SFT, the model $\pi^\text{SFT}$ is queried with inputs $x$ to sample pairs of responses $(y_1, y_2)\sim \pi^\text{SFT}(y \mid x)$. Human annotators then indicate a preference for one of the two outputs; the preferred and non-preferred responses under prompt $x$ are denoted $y_w\succ y_l \mid x$ for the winner and loser, respectively. These preference judgments are presumed to originate from an implicit, unobservable reward function $r^*(y, x)$. A variety of models exist for capturing such preferences, with the Bradley-Terry (BT) model \cite{bradley1952rankanalysis} being a widely adopted option. More general approaches, such as the Plackett-Luce models \citep{plackett1975analysis, luce2012individual}, can be utilized when working with multiple ranked answers. Within the BT framework, the probability of human preferences $p^*$ can be modeled as follows:

\begin{equation}\label{eq:bradley-terry}
    p^*(y_1\succ y_2 \mid x)=\frac{\exp\left(r^*(x, y_1)\right)}{\exp\left(r^*(x, y_1)\right) + \exp\left(r^*(x, y_2)\right)}.
\end{equation}

Given a fixed dataset of comparison tuples $\mathcal{D}=\bigl\{x^{(i)}, y_w^{(i)}, y_l^{(i)}\bigr\}_{i=1}^N$ drawn from $p^*$, we can introduce a reward model $r_{\phi}(x, y)$ parameterized by $\phi$, and estimate these parameters using the principle of maximum likelihood. Interpreting the problem as binary classification, we employ the following negative log-likelihood loss:

\begin{equation}\label{eq:reward_model}
    \mathcal{L}_R(r_{\phi}, \mathcal{D}) = -\mathbb{E}_{(x, y_w, y_l)\sim \mathcal{D}}\bigl[\log \sigma(r_{\phi}(x, y_w)- r_{\phi}(x, y_l))\bigr]
\end{equation}

where $\sigma$ denotes the logistic sigmoid. For language modeling tasks, $r_{\phi}(x, y)$ is typically initialized from the SFT model $\pi^\text{SFT}(y \mid x)$, with an extra linear projection layer appended to the terminal transformer output, which predicts a scalar reward value~\cite{ziegler2020finetuning}. To reduce the variance of the reward predictions, it is common practice to normalize the rewards across the data, enforcing $\mathbb{E}_{x,y\sim \mathcal{D}}\left[r_\phi(x, y)\right] = 0$ for all $x$.

\textbf{RL Fine-Tuning Phase}: In this stage, the model utilizes the learned reward as feedback to further adapt the language model. Consistent with previous studies~\citep{jaques2017sequence, jaques2020human}, the following optimization criterion is used:

\begin{equation}\label{eq:RL}
\max_{\pi_{\theta}}  \mathbb{E}_{x\sim \mathcal{D}, y\sim \pi_{\theta}(y \mid x)}\bigl[r_{\phi}(x, y)\bigr] - \beta\mathbb{D}_{\textrm{KL}}\bigl[\pi_{\theta}(y\mid x)\mid \mid \pi_\text{ref}(y\mid x)\bigr],
\end{equation}

with $\beta$ serving as a regularization parameter to control divergence from a reference policy $\pi_\text{ref}$, usually taken as the pre-trained SFT model $\pi^\text{SFT}$. The language model policy $\pi_\theta$ is normally initialized from $\pi^\text{SFT}$ as well. The regularization term plays a crucial role by constraining the policy to remain close to the region where the reward model provides accurate supervision and to avoid issues such as loss of diversity or collapse onto a narrow set of outputs. Because language output is inherently discrete, the objective is not directly differentiable; as such, reinforcement learning methods are generally required. The prevalent approach~\citep{ziegler2020finetuning, stiennon2022learning, bai2022training, ouyang2022training} formulates the reward as ${r(x, y) = r_{\phi}(x, y) -\beta (\log \pi_{\theta}(y\mid x) - \log \pi_\text{ref}(y\mid x))}$ and uses Proximal Policy Optimization (PPO)~\cite{schulman2017proximal} for maximization.

\section{Direct Preference Optimization}\label{sec:DPO}

Driven by the difficulties inherent in applying RL algorithms to large-scale tasks like language model fine-tuning, we seek a more direct alternative for optimizing policies from preference data. In contrast to standard RLHF pipelines that first train an explicit reward model and subsequently perform RL-based optimization, our methodology exploits a particular form of reward model parameterization that admits a closed-form solution for the associated optimal policy, thereby obviating the need for a traditional RL loop. Our central insight is that, by leveraging an explicit analytical connection between reward functions and their optimal induced policies, we can reformulate the loss on the reward model as an equivalent loss on the policy space itself. This change-of-variable perspective enables us to directly optimize the policy without fitting a separate reward model, while

\textbf{Deducing the DPO objective.} Beginning with the reinforcement learning objective established in previous literature, as given by Eq.~\ref{eq:RL}, and utilizing a general reward function $r$, it follows from earlier studies~\citep{peters2007reinforcement, peng2019advantage, korbak2022reinforcement, go2023aligning} that the optimizer for the KL-regularized reward maximization in Eq.~\ref{eq:RL} can be written as:

\begin{equation}\label{eq:op_policy}
    \pi_r(y\mid x) = \frac{1}{Z(x)}\pi_\text{ref}(y\mid x)\exp\left(\frac{1}{\beta}r(x, y)\right),
\end{equation}

where the normalization constant, or partition function, is $Z(x) =\sum_{y}\pi_\text{ref}(y\mid x)\exp\left(\frac{1}{\beta}r(x, y)\right)$. The full derivation is provided in Appendix \ref{app:derivation1}. Even substituting in an MLE estimator $r_{\phi}$ for the actual reward $r^*$, computing the partition function $Z(x)$ is computationally challenging~\citep{korbak2022reinforcement, go2023aligning}, which impedes the tractability of the resulting policy form in applications. Nevertheless, Eq.~\ref{eq:op_policy} can be reformulated to present the reward as a function of the optimal policy $\pi_r$, the baseline policy $\pi_\text{ref}$, and the (still implicit) partition function $Z(\cdot)$. By applying the logarithm to both sides of Eq.~\ref{eq:op_policy} and rearranging, we derive:

\begin{equation}\label{eq:main_eq}
    r(x,y) =\beta \log \frac{\pi_r(y\mid x)}{\pi_\text{ref}(y\mid x)} + \beta \log Z(x).
\end{equation}

Applying this transformation to the true reward $r^*$ and its associated optimal policy $\pi^*$, we note that the Bradley-Terry model only considers reward differences between two output candidates, namely, ${p^*(y_1 \succ y_2 \mid x) = \sigma(r^*(x, y_1) - r^*(x, y_2))}$. By substituting the representation from Eq.~\ref{eq:main_eq} into the Bradley-Terry preference equation, Eq.~\ref{eq:bradley-terry}, the terms involving the partition function are eliminated. This allows the probability of human preference to be reformulated solely in terms of the optimal policy $\pi^*$ and the reference $\pi_\text{ref}$. Hence, under the Bradley-Terry assumption, the optimal RLHF policy $\pi^*$ meets:

\begin{equation}\label{eq:objective}
    p^*(y_1\succ y_2 \mid x)=\frac{1}{1 + \exp\left(\beta \log \frac{\pi^*(y_2\mid x)}{\pi_\text{ref}(y_2\mid x)} - \beta \log \frac{\pi^*(y_1\mid x)}{\pi_\text{ref}(y_1\mid x)}\right)}
\end{equation}

A comprehensive derivation appears in Appendix~\ref{app:derivation2}. Although Eq.~\ref{eq:objective} is specific to the Bradley-Terry model, a similar approach yields formulations for broader preference models, including the Plackett-Luce family~\citep{plackett1975analysis, luce2012individual}, detailed in Appendix~\ref{app:plackett_luce_models}.

Having represented the likelihood of preference observations as a function of the optimal policy, one can now maximize the likelihood directly with respect to a parameterized policy $\pi_\theta$. By analogy to the reward-based modeling framework (see Eq.~\ref{eq:reward_model}), this leads to a policy learning objective as follows:

\begin{equation}\label{eq:optimum_model}
    \mathcal{L}_\text{DPO}(\pi_{\theta}; \pi_\text{ref}) = -\mathbb{E}_{(x, y_w, y_l)\sim \mathcal{D}}\left[\log \sigma \left(\beta \log \frac{\pi_{\theta}(y_w\mid x)}{\pi_\text{ref}(y_w\mid x)} - \beta \log \frac{\

In this approach, we model an implicit reward via an alternative parameterization, with its corresponding optimal policy given by $\pi_\theta$. Furthermore, because our method corresponds to fitting a reparameterized Bradley-Terry model, it inherits various theoretical guarantees, such as consistency provided certain conditions on the distribution of preference data are met \cite{bong2022generalized}. Section~\ref{sec:theory} elaborates further on the theoretical characteristics of DPO and compares them to related literature.

\textbf{What mechanism drives the DPO update?} To gain a deeper, mechanistic insight into DPO, one can examine the gradient of the loss function $\mathcal{L}_\text{DPO}$. The gradient with respect to the parameter vector $\theta$ is given by:

\begin{multline*}\label{eq:gradient}
    \nabla_\theta \mathcal{L}_\text{DPO}(\pi_\theta;\pi_\text{ref}) = \\ -\beta\mathbb{E}_{(x, y_w, y_l) \sim \mathcal{D}} \left[\underbrace{\sigma(\hat{r}_\theta(x, y_l) - \hat{r}_\theta (x, y_w))}_\text{greater influence when the reward ordering is inaccurate}\left[\underbrace{\nabla_\theta\log \pi(y_w \mid x)}_\text{push up probability of $y_w$} - \underbrace{\nabla_\theta\log\pi(y_l \mid x)}_\text{reduce probability of $y_l$}\right]\right],
\end{multline*}

where the implicit reward function $\hat{r}_\theta(x, y)$ is defined as $\beta \log \frac{\pi_\theta(y \mid x)}{\pi_\text{ref}(y \mid x)}$, determined by the current model $\pi_\theta$ relative to the reference $\pi_\text{ref}$ (see additional explanation in Section~\ref{sec:theory}). Conceptually, this gradient update increases the probability assigned to more preferred responses $y_w$, while suppressing those that are less preferred, $y_l$. Crucially, each example's contribution is modulated by the degree to which the implicit reward model $\hat{r}_\theta$ overvalues the less preferred completions, scaled by $\beta$—effectively emphasizing corrections where the model misorders preferences, modulated by the KL regularization strength. Our empirical findings underscore the necessity of this weighting: omitting the weighting factor can result in detrimental effects on the language model, as shown in Appendix Table~\ref{tab:unlikelihood_generations}.

\textbf{Overview of DPO.} The standard DPO workflow consists of two main steps: First, for each prompt $x$, generate two responses $y_1, y_2 \sim \pi_\text{ref}(\cdot \mid x)$ and obtain human preference labels to assemble the offline preference dataset $\mathcal{D} = \{x^{(i)}, y_w^{(i)}, y_l^{(i)}\}_{i=1}^N$. Second, train the language model $\pi_\theta$ by minimizing the loss $\mathcal{L}_\text{DPO}$ using the specified $\pi_\text{ref}$, $\mathcal{D}$, and parameter $\beta$.

In applied scenarios, practitioners often prefer leveraging existing, publicly available preference datasets rather than producing new completions and collecting fresh human annotations. When these datasets are built from samples drawn via $\pi^\text{SFT}$, we set the initial reference policy as $\pi_\text{ref} = \pi^\text{SFT}$ if possible. If $\pi^\text{SFT}$ is unavailable, we instead obtain $\pi_\text{ref}$ by fitting to the observed preferences, specifically maximizing the likelihood of preferred completions $(x, y_w)$, i.e., setting ${\pi_\text{ref} = \argmax_{\pi}\mathbb{E}_{x, y_w \sim \mathcal{D}}\left[\log \pi(y_w \mid x)\right]}$. This approach addresses the potential distribution mismatch between the true reference policy (which we lack direct access to) and the one utilized during DPO training. Additional implementation details and hyperparameter settings are described in Appendix~\ref{app:implementation}.

\section{Theoretical Insights into DPO}

This section develops a deeper conceptual understanding of the DPO framework, provides supporting theoretical arguments, and draws connections between the benefits of DPO and common challenges found in actor-critic based RLHF methods such as PPO~\cite{schulman2017proximal}.

\label{sec:theory}

\subsection{Your Language Model Is Secretly a Reward Model} DPO manages to avoid explicitly learning a reward function and sidesteps traditional reinforcement learning procedures by employing a single maximum likelihood training criterion. The objective function presented in Eq. \ref{eq:main_eq} can be interpreted as a Bradley-Terry model where the underlying reward is parameterized as $r^*(x, y) = \beta \log\frac{\pi^*_\theta(y \mid x)}{\pi_\text{ref}(y \mid x)}$. Training our model $\pi_{\theta}$ in this way is mathematically analogous to reward model optimization as shown in Eq. \ref{eq:reward_model}, given a suitable change of variables. In the discussion below, we formalize the theoretical justification for this reparameterization, demonstrate that it imposes no limitations on the class of expressible reward models, and establish that this method allows for perfect recovery of the optimal policy. We start by formalizing a notion of equivalence between reward functions.

\begin{definition}
Two reward functions $r(x, y)$ and $r'(x, y)$ are defined as equivalent if and only if there exists a function $f$ such that $r(x, y) - r'(x, y) = f(x)$.
\end{definition}

It is straightforward to verify that this forms an equivalence relation, which divides the collection of reward functions into equivalence classes. With this setup, we establish two key lemmas:

\begin{lemma}\label{lemma:same_prefrence}
Within the Plackett-Luce model, and in particular the Bradley-Terry framework, any pair of reward functions belonging to the same equivalence class induces an identical preference distribution.
\end{lemma}

\begin{lemma}\label{lemma:same_policy}
    Any two reward functions in the same equivalence class yield the same optimal policy in the corresponding constrained reinforcement learning (RL) problem.
\end{lemma}

The arguments are elementary and can be found in Appendix~\ref{app:lemma1}. The first lemma highlights a classical identification issue inherent to the Plackett-Luce model family \cite{plackett1975analysis}. As a result, one typically needs to enforce additional constraints for identifiability to ensure that MLE solutions derived from Eq.~\ref{eq:reward_model} are meaningful \cite{bong2022generalized}. The second lemma reveals that every reward function in a particular equivalence class leads to the same optimal policy. Therefore, our main objective reduces to identifying any representative reward function from the optimal class. We formalize this in the following theorem, which is proved in Appendix~\ref{app:thm1}:

\begin{theorem}\label{thm:main}
    Subject to mild conditions, all reward function equivalence classes compatible with the Plackett-Luce (or Bradley-Terry) models can be expressed using the reparameterization ${r(x, y) = \beta \log \frac{\pi(y\mid x)}{\pi_\text{ref}(y\mid x)}}$ for some probability model $\pi(y\mid x)$ and a fixed reference model $\pi_\text{ref}(y\mid x)$.
\end{theorem}

\begin{sproof}
    Given any reward function $r(x, y)$, one can construct the associated optimal policy model $\pi_r(y \mid x)$ as defined in Eq.~\ref{eq:op_policy}. We aim to demonstrate that there exists an equivalent reward function (belonging to the same class as $r$) that admits the aforementioned reparameterized form. Consider the projection operator $f$ defined as
\begin{equation}
    f(r; \pi_\text{ref}, \beta)(x, y) = r(x, y) - \beta \log \sum_{y} \pi_\text{ref}(y \mid x) \exp\left(\frac{1}{\beta} r(x, y)\right)
\end{equation}
This mapping effectively adjusts $r$ by subtracting the log-partition function (a normalization term) computed under $\pi_\text{ref}$, depending only on $x$. Consequently, $f(r; \pi_\text{ref}, \beta)(x, y)$ remains in the same equivalence class as $r(x, y)$. By substituting $r$ with the right-hand side of Eq.~\ref{eq:main_eq} (which universally applies to any reward function), we find $f(r; \pi_\text{ref}, \beta)(x, y) = \beta \log \frac{\pi_r(y\mid x)}{\pi_\text{ref}(y\mid x)}$. Therefore, this construction yields a representative from the equivalence class of $r$ in the required form, ensuring no loss of generality in our choice of reward parameterization.
\end{sproof}

An alternative interpretation of Theorem~\ref{thm:main} is that it identifies the particular reward function chosen by the DPO reparameterization within each equivalence class—specifically, the one satisfying:

\begin{equation}\label{eq:lag_p}
     \sum_{y}\underbrace{\pi_\text{ref}(y\mid x)\exp\left(\frac{1}{\beta}r(x, y)\right)}_{=\pi(y\mid x)\text{, using Thm.~\ref{thm:main} reparam.}} = 1,
\end{equation}

This means that $\pi(y\mid x)$ indeed constitutes a valid probability distribution: all entries are nonnegative, and the sum totals to one. Furthermore, referring to Eq.~\ref{eq:op_policy}, it becomes evident that Eq.~\ref{eq:lag_p} describes the partition function corresponding to the optimal policy that is determined by the reward function $r(x, y)$. The principal insight leveraged by the DPO algorithm lies in enforcing specific constraints on the otherwise under-constrained Plackett-Luce (and, more specifically, Bradley-Terry) class of preference models. By doing so, it maintains the set of realizable reward models, while ensuring that the optimal policy characterized in Eq.~\ref{eq:op_policy} is computationally accessible for any prompt $x$.

\subsection{Instability of Actor-Critic Algorithms}
Our analytical framework also provides a lens to examine and explain the instabilities that often arise in conventional actor-critic algorithms utilized in RLHF procedures, such as PPO. Here, we follow the established RLHF workflow and concentrate on the RL fine-tuning process as outlined in Section \ref{section:prelims}. Connections can be drawn to the control as inference paradigm \cite{levine2018reinforcement}, specifically as applied to the constrained RL formulation introduced in \ref{eq:RL}. Considering a parameterized policy $\pi_{\theta}(y\mid x)$, the objective is to minimize $\mathbb{D}_{\text{KL}}[\pi_{\theta}(y|x) \mid \mid \pi^*(y\mid x)]$, where $\pi^*$ denotes the optimal policy derived from Eq.~\ref{eq:optimum_model}, based on the reward $r_{\phi}(y, x)$. Through algebraic manipulation, this results in the following maximization objective:

\begin{equation}\label{eq:AC}
    \max_{\pi_{\theta}}\mathbb{E}_{\pi_{\theta}(y\mid x)}\bigg[\underbrace{r_{\phi}(x, y) -\beta\log\sum_{y}\pi_\text{ref}(y\mid x)\exp\left(\frac{1}{\beta}r_{\phi}(x, y)\right)}_{f(r_{\phi}, \pi_\text{ref}, \beta)} - \underbrace{\beta\log\frac{\pi_{\theta}(y\mid x)}{\pi_\text{ref}(y\mid x)}}_{\text{KL}}\bigg]
\end{equation}

This objective is identical to that targeted in previous studies 
\citep{ziegler2020finetuning, stiennon2022learning, bai2022training, ouyang2022training}, where the DPO-equivalent reward is utilized for the reward category $r_{\phi}$. Within this framework, the normalization component in $f(r_{\phi}, \pi_\text{ref}, \beta)$ can be viewed as representing the soft value function associated with the reference policy $\pi_\text{ref}$. Although the normalization term does not modify the optimal solution, omitting it can result in high-variance policy gradients, which destabilizes the learning process. One could address this normalization by employing a value function approximator, yet optimizing such an estimator poses significant challenges. Alternatively, earlier works have normalized rewards by comparing to a baseline generated by a human completion, effectively yielding a single-sample Monte Carlo estimate of the normalization constant. Notably, with the DPO reparameterization, the resulting reward function inherently avoids the need for any baseline normalization.

\section{Experiments}
This section provides an empirical assessment of DPO's effectiveness for training policies directly based on preference data. We begin by investigating a controlled text-generation scenario, exploring how DPO balances reward maximization and KL-divergence minimization with the reference policy, as compared to prevalent algorithms like PPO for preference learning. Subsequently, we extend our evaluation to larger models and more challenging RLHF tasks, such as summarization and dialogue. Our results indicate that DPO, requiring minimal hyperparameter adjustment, consistently matches or surpasses established baselines—including RLHF with PPO and selecting the best out of $N$ sampled rollouts according to a learned reward. Prior to detailing these findings, we outline the experimental methodology; further information can be found in Appendix~\ref{app:exp_details}.

\textbf{Tasks.} We investigate three distinct open-domain text generation tasks in our experimental setup. Across all settings, the algorithms are trained to optimize a policy using a dataset of preferences $\mathcal{D}=\bigl\{x^{(i)}, y_w^{(i)}, y_l^{(i)}\bigr\}_{i=1}^N$. In the \textbf{controlled sentiment generation} task, $x$ serves as an initial segment of a movie review extracted from the IMDb dataset \cite{maas-EtAl:2011:ACL-HLT2011}, and the objective for the policy is to produce a continuation $y$ that conveys positive sentiment. To facilitate systematic evaluation in this task, preference pairs over outputs are \textit{generated} via a pre-trained sentiment classifier, enforcing the criterion that $p(\text{positive}\mid x,y_w)>p(\text{positive}\mid x,y_l)$. For the SFT baseline, we apply additional training to GPT-2-large until it converges, using the training partition of IMDb reviews (additional implementation specifics are found in App~\ref{app:sentiment_details}). In the \textbf{summarization} scenario, $x$ consists of a Reddit forum post, and the policy must produce a concise summary $y$ encapsulating the central points. Consistent with previous studies, the Reddit TL;DR summarization dataset \citep{volske-etal-2017-tl} and the human preferences collected by \citeauthor{stiennon2022learning} are employed. An SFT model fine-tuned on human-authored forum summaries\footnote{\url{https://huggingface.co/CarperAI/openai_summarize_tldr_sft}} is utilized, with RLHF implemented through the TRLX \citep{leandro_von_werra_2023_7790115} toolkit. Notably, the human preference annotations originate from work by \citeauthor{stiennon2022learning}, using generations from a similar yet distinct SFT-trained model. For the \textbf{single-turn dialogue} task, $x$ represents a user query, ranging from inquiries about astrophysics to requests for advice. Here, the policy is required to generate a response $y$ that is both useful and engaging; we utilize the Anthropic Helpful and Harmless dialogue dataset \citep{bai2022training}, which encompasses 170,000 conversational exchanges between a human user and an AI assistant. Each dialogue concludes with two AI-generated responses and an associated human preference label. Unlike previous settings, a pre-existing SFT model is unavailable for this task, so we instead fine-tune a generic language model using only the human-preferred continuations to construct the SFT model.

\textbf{Evaluation.} We employ two distinct evaluation methodologies in our experimental analysis. For the controlled sentiment generation scenario, where the true reward function (provided by a sentiment classifier) is accessible, we quantify each algorithm’s performance by mapping the trade-off curve between attained reward and KL-divergence relative to the reference policy. This reward-divergence frontier is computable owing to our access to the ground-truth reward. However, this is not representative of real-world applications, where the underlying reward function remains unknown. In these situations, we assess algorithms using their \textit{win rate} when compared to a baseline policy. Here, we utilize GPT-4 as an automated proxy for human judgments—evaluating the quality of summaries and the helpfulness of responses in the summarization and single-turn dialogue tasks, respectively. For summarization, baseline comparisons employ the reference summaries from the test set; for dialogue, the baseline is the test set’s preferred response. Although previous works indicate that LMs are often superior to traditional metrics as automated evaluators \citep{Chen2023ExploringTU}, we supplement this with a human study, detailed in Sec.~\ref{sec:human-judgments}, to substantiate our reliance on GPT-4 for evaluation. Our results show that GPT-4 assessments exhibit a strong alignment with human judgment, with levels of agreement between humans and GPT-4 equalling or exceeding inter-annotator consensus among humans.

\textbf{Methods.} Beyond DPO, we benchmark several established strategies for training language models to reflect human preference. For instance, we investigate zero-shot prompting using \textbf{GPT-J} \citep{gpt-j} for the summarization task and apply 2-shot prompting with \textbf{Pythia-2.8B} \citep{biderman2023pythia} for dialogue. Our evaluation also encompasses the \textbf{SFT} approach, alongside \textbf{Preferred-FT}, a variant fine-tuned with supervised learning on the selected completion $y_w$. The source of $y_w$ is either the SFT model (for controlled sentiment and summarization) or a generic language model (for single-turn dialogue). We further examine a pseudo-supervised technique called \textbf{Unlikelihood}~\citep{welleck2019neural}, which explicitly boosts the probability of $y_w$ while reducing that of $y_l$, incorporating an optional scaling parameter $\alpha \in [0,1]$ to adjust the contribution of the unlikelihood component. The comparison set also includes \textbf{PPO} \citep{schulman2017proximal}, leveraging a reward function inferred from preference annotations, and \textbf{PPO-GT}, which operates as an oracle trained directly on the known reward function in the controlled sentiment context. For our sentiment tasks, we deploy two versions of PPO-GT: a standard, out-of-the-box implementation \cite{leandro_von_werra_2023_7790115}, and a customized variant that normalizes reward signals and undergoes additional hyperparameter optimization for enhanced outcomes (these modifications are similarly applied when running PPO with learned rewards). Additionally, we introduce the \textbf{Best of $N$} method as a baseline: this involves generating $N$ candidate outputs from the SFT model (or Preferred-FT in the dialogue setting) and selecting the highest-rewarded output, as per a reward model trained on preference data. While this baseline is often competitive and allows reward model quality to be separated from PPO’s optimization process, its computational demands are prohibitive for even moderate values of $N$, as it necessitates $N$ samples for every test-time query.

\subsection{How well can DPO optimize the RLHF objective?}

The standard KL-constrained reward maximization objective underlying most RLHF methodologies is designed to strike a balance between maximizing the reward and limiting the policy's divergence from the reference policy. As such, a meaningful evaluation of different algorithms requires considering both the reward obtained and the magnitude of the KL divergence; pursuing higher rewards at the cost of significantly increased KL is often not optimal. Figure~\ref{fig:frontier-tldr-main} presents the reward-KL tradeoff curves for several algorithms in the sentiment analysis context. For each algorithm, we conduct several training runs, varying the hyperparameters governing policy conservativeness for each trial (target KL $\in\{3,6,9,12\}$ for PPO, $\beta \in \{0.05,0.1,1,5\}$, $\alpha\in\{0.05,0.1,0.5,1\}$ for unlikelihood, and different random seeds for preferred-FT), resulting in a total of 22 experimental runs. After every 100 optimization steps up to convergence, we evaluate each resulting policy over a suite of held-out prompts, measuring both the mean reward according to the actual reward function and the average sequence-level KL divergence\footnote{This represents the aggregate of the KL divergence across all timesteps in the sequence.} relative to the reference policy, expressed as $\text{KL}\left(\pi\mid \mid \pi_\text{ref}\right)$. Our analysis indicates that DPO yields a substantially superior Pareto frontier, achieving the greatest rewards while maintaining relatively low KL values. This observation is especially significant for several reasons. Notably, although DPO and PPO optimize an equivalent objective, DPO demonstrates considerably greater efficiency, with DPO’s reward-to-KL efficiency consistently outperforming that of PPO. Furthermore, DPO outperforms PPO on the reward-KL tradeoff curve \emph{even in scenarios where PPO has access to exact reward information} (PPO-GT).

\subsection{Can DPO scale to real preference datasets?}
\label{sec:dpo-real-datasets}
We further investigate how well DPO performs when fine-tuned on tasks in summarization and single-turn dialogue settings. Notably, commonly used automatic evaluation metrics in summarization, such as ROUGE, often have a weak correlation with actual human judgments~\citep{stiennon2022learning}. Previous studies indicate that leveraging PPO for language model fine-tuning based on human preferences yields more effective summaries. To benchmark different techniques, we sample model outputs from the test portion of the TL;DR summarization dataset and calculate their mean win rate compared to references in the test data. Outputs for each approach are sampled at a range of temperatures between 0.0 and 1.0; the win rates corresponding to these sampling regimes are depicted in Figure~\ref{fig:frontier-tldr-main} (right). All methods, including DPO, PPO, and Preferred-FT, start from the same GPT-J SFT checkpoint\footnote{\url{https://huggingface.co/CarperAI/openai_summarize_tldr_sft}} for fine-tuning. Our experiments show that DPO attains a win rate close to 61\% when sampling at temperature 0.0, outperforming PPO, which achieves a roughly 57\% win rate at its best sampling temperature of 0.0. Additionally, DPO surpasses the best-of-$N$ baseline in terms of maximum win rate. It is important to mention that no substantial effort was made to optimize DPO's $\beta$ hyperparameter, suggesting these results might be conservative estimates of DPO's capabilities. Furthermore, DPO demonstrates much greater resilience to changes in sampling temperature than PPO, whose performance can drop to that of the baseline GPT-J at higher temperatures. Preferred-FT yields minimal improvement relative to the initial SFT model. In addition, Section~\ref{sec:human-judgments} presents a direct comparison between DPO and PPO using human assessments, where DPO outputs at temperature 0.25 are preferred in 58\% of cases compared to PPO outputs at temperature 0.

For single-turn dialogue evaluation, we assess the various methods using the subset of the Anthropic HH test split \citep{bai2022training} containing only a single round of interaction between a human and the assistant. The GPT-4 assessments rely on preferred completions from the test data as references, calculating win rates across different methods. Lacking an established SFT baseline for this particular task, our approach begins with a pre-trained Pythia-2.8B; we apply Preferred-FT to train a reference model on the selected completions to ensure model outputs remain in-distribution, followed by DPO training. We include, for comparison, both the best completion selected from 128 Preferred-FT samples (noting that the Best of $N$ baseline reaches a plateau at $N=128$, as detailed in Appendix Figure~\ref{fig:best-of-n}) and a two-shot prompted Pythia-2.8B base variant, observing that DPO meets or exceeds their results at optimal temperatures. Additionally, an RLHF model trained with PPO on the Anthropic HH dataset \footnote{\url{https://huggingface.co/reciprocate/ppo_hh_pythia-6B}}, sourced from a recognized repository \footnote{\url{https://github.com/CarperAI/trlx/tree/main/examples/hh}}, is evaluated; however, no prompt or sampling temperature yields better outcomes than the unmodified Pythia-2.8B. Given both our TL;DR findings and the fact that PPO and Best of 128 optimize the same reward structure, we treat Best of 128 as an approximate stand-in for PPO performance. In sum, DPO is distinguished as the only resource-efficient technique that consistently surpasses preferred completions on the Anthropic HH dataset and achieves performance on par with or superior to the computationally intensive Best of 128 method. Lastly, Figure~\ref{fig:dialogue-main} demonstrates that DPO achieves optimal performance in relatively few training steps.

\subsection{Generalization to a new input distribution}

To further investigate how PPO and DPO handle shifts in input distribution, we assessed the policies obtained from our Reddit TL;DR summarization study on a distinct dataset: the test split of CNN/DailyMail news articles \citep{nallapati-etal-2016-abstractive}. We used the best-performing sampling temperatures from the TL;DR task (0 and 0.25) for evaluation. The corresponding outcomes can be found in Table~\ref{tab:ood}. To measure performance, we calculated GPT-4's win rate against the reference summaries in the dataset, employing the same GPT-4 (C) evaluation prompt from Reddit TL;DR, modified only by replacing ``forum post'' with ``news article''. Across this new data distribution, DPO still significantly surpasses PPO in performance. These findings offer preliminary support that DPO can achieve generalization performance comparable to PPO, despite not leveraging the supplementary unlabeled Reddit TL;DR prompts that PPO incorporates.

\subsection{Validating GPT-4 judgments with human judgments}
\label{sec:human-judgments}

We ran a human evaluation to assess the trustworthiness of GPT-4's comparative judgments using results from the TL;DR summarization task with two variants of GPT-4 prompts. The \textbf{GPT-4 (S)} (simple) prompt inquires solely about which summary better captures the critical information from the post, while the \textbf{GPT-4 (C)} (concise) prompt additionally includes a requirement for conciseness; the latter is used because GPT-4, under the (S) prompt, tends to select longer and more repetitive outputs than human raters do. Full prompt text is available in Appendix~\ref{app:prompts}. To ensure diverse output quality, we evaluated three methods: DPO at temperature 0.25 (highest-performing), PPO at temperature 1.0 (lowest), and SFT at temperature 0.25 (intermediate), each compared against PPO's greedy decoding baseline (its optimal temperature). Our analysis shows that, regardless of the prompt used, GPT-4’s selections typically align with human preferences at rates similar to inter-human agreement, indicating GPT-4 can serve as an adequate surrogate for human evaluation (we obtained multiple human ratings only for the DPO and PPO-1 pairings, due to the small pool of raters). Notably, the \textbf{GPT-4 (C)} prompt delivers win rates that more closely match human preferences, which motivates its use for the principal comparisons in Section~\ref{sec:dpo-real-datasets}. Further methodological details, including information about the rater-facing interface and volunteer participants, can be found in Appendix~\ref{app:human-study}.

\section{Discussion}
Preference-based learning provides an effective and adaptable approach for training robust and well-aligned language models. In this work, we presented DPO, an intuitive training scheme designed to leverage preference data for language model optimization, circumventing the complexities of reinforcement learning. Instead of reshaping the preference learning task into a conventional RL context to accommodate standard RL techniques, DPO establishes a direct correspondence between language model policy outputs and reward functions, allowing language models to align with human preferences using a straightforward cross-entropy objective. This process eliminates the need for reinforcement learning, all while maintaining full generality. DPO matches or surpasses the performance of popular RLHF approaches, such as those employing PPO, and achieves these results with minimal hyperparameter adjustment, thereby substantially lowering the barriers to preference-based language model development.

\textbf{Limitations \& Future Work.} Our findings open avenues for further exploration. For instance, how robustly does a policy trained via DPO generalize to out-of-distribution data compared to models optimized with explicit reward functions? Preliminary evidence indicates that DPO policies exhibit generalization on par with those derived from PPO, yet a broader evaluation is necessary. Another question is whether techniques like self-labeling based on DPO-trained policies can successfully exploit unlabeled input prompts. Additionally, it remains unclear how reward over-optimization appears in the DPO paradigm; the marginal performance drop observed in Figure~\ref{fig:dialogue-main}-right might be indicative of such effects. Moreover, although our experiments involve models up to 6B parameters, scaling DPO to support significantly larger, state-of-the-art models represents a promising line of inquiry. On the evaluation front, we observed that GPT-4-based win rates depend on the prompt formulation, suggesting a need for further research into eliciting reliable automated evaluations. Beyond training language models from human preferences, DPO also holds promise for extending preference learning to other generative modeling domains.